\section{Адаптивное и робастное управление}
\label{sec:adaptive}

\subsection{Адаптация усиления при неизвестном параметре}
\lecture{L6}{15--17}

Адаптивный регулятор изменяет свои параметры во времени, используя измерения
объекта. В примере лекции требуется стабилизировать скалярную систему
\[
 \dot x=ax+u,\qquad a=-a_0+\Delta a,\qquad a_0>0,
\]
где $\Delta a$ --- неизвестное постоянное число. Если бы $a$ было известно,
можно было бы сразу выбрать достаточно сильную отрицательную обратную связь.
Адаптивный закон позволяет обойтись без известной верхней границы $\Delta a$:
\begin{equation}
 u=-c(t)x,\qquad \dot c=\gamma x^2,\qquad \gamma>0.
 \label{eq:adaptive-law}
\end{equation}
Здесь $c$ --- состояние регулятора, а $\gamma$ --- коэффициент скорости
адаптации. Это не входной коэффициент $\gamma(x)$ из линеаризации обратной
связью: обозначение используется в другом локальном контексте лекции.

Подстановка закона даёт
\[
 \dot x=-a_0x-(c-\Delta a)x.
\]
Вводится функция Ляпунова для состояния и отклонения усиления:
\[
 V(x,c)=\frac{x^2}{2}+\frac{(c-\Delta a)^2}{2\gamma}.
\]
Параметр $\Delta a$ входит в доказательство, но не требуется для вычисления
управления или $\dot c$. Поскольку $\Delta a$ постоянно,
\begin{align*}
 \dot V&=x[-a_0x-(c-\Delta a)x]
       +\frac{c-\Delta a}{\gamma}\dot c\\
 &=-a_0x^2+\frac{c-\Delta a}{\gamma}(\dot c-\gamma x^2)
 =-a_0x^2\le0.
\end{align*}
Правило адаптации выбрано именно для сокращения неизвестного смешанного
слагаемого. Если $x\ne0$, усиление $c$ увеличивается. Когда ошибка уменьшается,
скорость адаптации также уменьшается. В иллюстрации лекции используется
частный случай $\gamma=1$, то есть $\dot c=x^2$.

\addition Из $V(t)\le V(0)$ следуют ограниченность $x$ и $c$ и полнота
решения. В множестве $\dot V=0$ имеем $x=0$; оно инвариантно при любом
постоянном $c$. По принципу ЛаСалля $x(t)\to0$. В силу $\dot c\ge0$ и
ограниченности $c$ имеет конечный предел. Но из доказательства не следует
$c(t)\to\Delta a$: на $x=0$ адаптация останавливается при любом значении
усиления. Стабилизация состояния и идентификация неизвестного параметра
не являются одним и тем же результатом.

Также неверно утверждать, что $\dot V<0$ для любого ненулевого
$(x,c-\Delta a)$. При $x=0$, $c\ne\Delta a$ производная равна нулю.
Поэтому строгий вывод относится к сходимости $x$, а не к асимптотической
устойчивости единственной точки $(0,\Delta a)$ расширенной системы.

\subsection{Адаптивный закон с релейной составляющей}
\lecture{L6}{18--19}

Лекция дополняет управление постоянным релейным усилением:
\[
 u=-cx-d\sgn x,\qquad d>0,\qquad \dot c=\gamma x^2.
\]
Для прежней функции Ляпунова получается
\[
 \dot V=-a_0x^2-dx\sgn x=-a_0x^2-d|x|\le0.
\]
Квадратичное слагаемое отвечает за обычное затухание, дополнительное
$-d|x|$ усиливает притяжение при малых ненулевых $x$.
Закон опять не использует $\Delta a$. При $x=0$ разрыв трактуется
как в разделе о скользящем режиме.

\addition Нельзя автоматически подставить полную $V(x,c)$ в формулу
конечного времени с $\dot V\le-k\sqrt V$: величина $V$ включает ошибку
усиления, которая может остаться ненулевой после $x=0$.
Наличие в $\dot V$ слагаемого $-|x|$ не является само по себе доказательством
конечного времени обращения всей расширенной функции $V$ в нуль.

\subsection{Робастное усиление при известном интервале параметра}
\lecture{L6}{20--21}

Робастный регулятор использует известное множество возможных параметров
и обеспечивает результат для любого элемента этого множества.
Для $\dot x=ax+u$ пусть
\[
 \underline a\le a\le\overline a.
\]
Выбираем постоянное усиление
\[
 u=-kx,\qquad k>\overline a.
\]
Тогда $\dot x=-(k-a)x$, а решение удовлетворяет равномерной оценке
\[
 |x(t)|=|x(0)|e^{-(k-a)t}
 \le|x(0)|e^{-(k-\overline a)t}.
\]
Доказана глобальная экспоненциальная устойчивость для всего интервала.
В примере лекции указан закон $u=-3x$; он подходит, когда известная верхняя
граница $a$ строго меньше $3$. Для произвольного неизвестного и ничем
не ограниченного $a$ такое фиксированное усиление гарантии не даёт.

В отличие от адаптации здесь нет собственного состояния регулятора:
усиление выбирается один раз по худшему допустимому случаю. Этот выбор
может быть консервативным, зато доказательство сразу охватывает весь
заданный диапазон.

\subsection{Робастное подавление ограниченного возмущения}
\lecture{L6}{22--25}

Теперь объект имеет вид
\[
 \dot x=ax+u+f(t),\qquad
 a\in[\underline a,\overline a],\qquad |f(t)|\le\alpha.
\]
При одном $u=-kx$ постоянное $f$ оставило бы установившееся смещение.
В лекции выбирается
\begin{equation}
 u=-kx-k_0\sgn x,\qquad k>\overline a,\quad k_0>\alpha.
 \label{eq:robust-scalar}
\end{equation}
Для $V=x^2/2$ имеем
\begin{align*}
 \dot V&=-(k-a)x^2-k_0|x|+xf(t)\\
 &\le-(k-\overline a)x^2-(k_0-\alpha)|x|.
\end{align*}
Неизвестное возмущение оценивается по модулю. Строгий запас
$k_0-\alpha>0$ сохраняет притяжение даже при наихудшем знаке $f$.
Поскольку возмущение входит в тот же канал, что и управление, релейная
часть способна поддерживать идеальное скольжение в нуле.

\addition Для $x\ne0$ из этой же оценки следует
\[
 \frac{d}{dt}|x|\le-(k-\overline a)|x|-(k_0-\alpha)
 \le-(k_0-\alpha).
\]
Следовательно, в смысле Филиппова нуль достигается за
$t_*\le|x(0)|/(k_0-\alpha)$.
В численной иллюстрации лекции
\[
 u=-3x-4\sgn x,\qquad f(t)=2\sin50t.
\]
Здесь $\alpha=2$, запас реле равен $2$, а условие линейной части ---
$\overline a<3$. Возмущение быстро меняется, но для оценки важна его
амплитуда, а не возможность вычислить его мгновенное значение.
Идеальная оценка времени даёт $t_*\le|x(0)|/2$.

Если заменить знак насыщением, компенсация в самом нуле становится
однозначной и нулевой; произвольное ненулевое $f$ уже может вывести систему
из нуля. Тогда необходимо заново получить оценку окрестности или доказать
точную сходимость при дополнительных свойствах неопределённости.
Именно это различие используется при разборе двух непрерывных регуляторов
работы 5.

\subsection{Как сопоставить подходы на защите}

\begin{center}
\begin{tabularx}{\textwidth}{@{}lYY@{}}
\toprule
Подход & Что заранее известно & Как используется неопределённость\\
\midrule
Точная компенсация & Нелинейности и коэффициент входа & Вычисляются непосредственно в законе\\
Робастное усиление & Гарантированные границы & Выбирается запас для худшего случая\\
Скользящий режим & Знак входа и оценка нормированной ошибки & Релейная часть превосходит эту оценку\\
Адаптация & Структура модели и доступные измерения & Усиление изменяется по измеряемой ошибке\\
\bottomrule
\end{tabularx}
\end{center}
Эти подходы могут сочетаться: достаточно точная номинальная компенсация
может уменьшить релейную часть, адаптация изменяет усиление,
насыщение делает сигнал непрерывным.
Но для каждой полученной комбинации требуется своё доказательство;
свойства одного составляющего метода нельзя переносить на весь контур
без проверки.
